%========================================================================================
%  CAPITOLO 4 — INFERENZA BAYESIANA E STRUMENTI DI ANALISI
%  Lunghezza indicativa: 5-6 pagine.
%  Fonti: Thrane & Talbot (2019), Taylor (2021) cap. 5-7, tesi di Irene cap. 3 (impostazione).
%========================================================================================

\chapter{Inferenza bayesiana e strumenti di analisi} \label{chpt:inferenza}
\sloppy

\section{Teorema di Bayes, likelihood e prior} \label{sec:bayes}
% - Posterior = likelihood x prior / evidenza. Marginalizzazione. Intervalli di credibilità.

\section{La likelihood per il pulsar timing} \label{sec:likelihood_PTA}
% - Residui gaussiani: r = M eps + F a + n.
% - Likelihood marginalizzata sui coefficienti: ln L = -1/2 r^T C^-1 r - 1/2 ln det(2 pi C),
%   C = N + T B T^T (formula di Woodbury per l'inversione).

\section{Campionamento: Markov Chain Monte Carlo} \label{sec:MCMC}
% - Metropolis-Hastings, PTMCMC (parallel tempering, salti adattivi).
% - Impostazioni usate: Niter = 1e6, thin = 10, burn-in 25%, ~1e5 campioni.
% - Convergenza: tempo di autocorrelazione (emcee), corner plot.

\section{Enterprise e Discovery} \label{sec:enterprise_discovery}
% - Enterprise: standard della comunità (Python/NumPy).
% - Discovery: nuova implementazione in JAX, differenziabile, GPU, compilazione JIT.
% - Perché è importante validarla -> capitolo 5.
% - Workstation di dipartimento (GPU) per i run lunghi.
